\chapter*{Introduction générale}
\addcontentsline{toc}{chapter}{Introduction générale}
\markboth{Introduction générale}{}

La transition écologique constitue aujourd'hui un enjeu majeur pour les sociétés modernes. Face à l'augmentation continue de la production des déchets, à l'épuisement des ressources naturelles et aux conséquences environnementales des modes de consommation actuels, l'adoption de pratiques responsables devient indispensable. Le tri sélectif et la valorisation des déchets représentent, dans ce contexte, des leviers essentiels pour réduire l'enfouissement, favoriser le recyclage et contribuer à la préservation de l'environnement.

Cependant, l'efficacité des politiques de gestion des déchets ne dépend pas uniquement de la disponibilité des infrastructures de collecte. Elle repose également sur l'implication active des citoyens, leur sensibilisation aux consignes de tri et leur capacité à accéder rapidement à des informations fiables. Dans la pratique, plusieurs difficultés persistent : l'identification des catégories de déchets, la méconnaissance des règles de tri, l'absence d'informations centralisées sur les points de collecte et le manque de mécanismes susceptibles de maintenir l'engagement des utilisateurs dans la durée.

L'essor des technologies numériques offre des opportunités concrètes pour répondre à ces problématiques. Les applications mobiles et web permettent aujourd'hui de diffuser des informations personnalisées, de faciliter l'accès aux services de proximité et de proposer des expériences interactives favorisant l'adoption de nouveaux comportements. Elles peuvent également intégrer des mécanismes de gamification, de suivi et de participation communautaire afin de transformer les gestes écologiques en actions visibles, mesurables et valorisées.

C'est dans cette perspective que s'inscrit le présent projet de fin d'études, réalisé au sein d'\textbf{ASSURANCES MAGHREBIA}. Au-delà de son activité principale dans le domaine de l'assurance des biens et des personnes, l'entreprise s'inscrit dans une démarche responsable intégrant les principes du développement durable et de l'innovation numérique. Cette orientation a conduit à la conception d'\textit{EcoRewind}, une plateforme numérique dédiée à la sensibilisation, à l'accompagnement et à la valorisation du tri sélectif des déchets en Tunisie.

\textit{EcoRewind} est une application mobile et web conçue selon une approche centrée sur l'utilisateur. Elle vise à proposer une expérience simple, accessible et motivante, permettant aux citoyens d'adopter plus facilement des comportements écoresponsables. La plateforme met à disposition une carte interactive permettant de localiser les points de collecte selon les types de déchets, un guide de tri ainsi qu'un scanner destiné à aider l'utilisateur dans l'identification des déchets. Elle propose également un fil d'actualité communautaire favorisant le partage d'initiatives écologiques, de conseils et de témoignages entre les utilisateurs.

Afin d'encourager l'engagement durable des citoyens, la solution intègre des mécanismes de gamification fondés sur l'accumulation de points, l'obtention de badges et le suivi de la progression individuelle. Ces fonctionnalités permettent de valoriser les actions réalisées et de rendre l'impact environnemental de chaque utilisateur plus concret, notamment à travers des indicateurs relatifs aux déchets triés et au CO\textsubscript{2} évité. La plateforme comprend également un espace d'administration dédié à la gestion des utilisateurs, des publications, des témoignages et des points de collecte.

Dans le souci de garantir la qualité et la sécurité des échanges au sein de la communauté, \textit{EcoRewind} intègre un mécanisme de modération des contenus. Celui-ci associe des règles métier à des techniques d'intelligence artificielle pour l'analyse des textes et des images, tout en prévoyant une validation humaine des contenus nécessitant une vérification. Cette approche contribue à assurer la pertinence des publications et à maintenir un environnement d'interaction fiable et sécurisé.

Dans le cadre de ce projet, un prototype fonctionnel de la plateforme \textit{EcoRewind} a été conçu et réalisé afin de valider les principaux choix fonctionnels, ergonomiques et techniques. Ce prototype met en œuvre les parcours essentiels de l'utilisateur, notamment la consultation des points de collecte à travers une carte interactive, l'accès aux informations relatives au tri, la participation à l'espace communautaire ainsi que le suivi de la progression individuelle. Il constitue une première concrétisation de la solution proposée et permet d'évaluer la cohérence des interfaces, la pertinence des fonctionnalités développées ainsi que l'intégration des différents composants techniques. Cette étape de prototypage représente une base essentielle pour recueillir les retours des utilisateurs, identifier les axes d'amélioration et orienter les évolutions futures de la plateforme.

La réalisation de cette solution s'appuie sur une démarche agile fondée sur le cadre \textbf{Scrum}. Cette méthodologie permet d'organiser le travail de manière itérative et incrémentale, en découpant le projet en sprints et en priorisant les fonctionnalités selon leur valeur. Elle favorise également une amélioration continue de la solution grâce aux retours recueillis tout au long du processus de développement.

Sur le plan technique, \textit{EcoRewind} est développée avec le framework \textbf{Flutter}, qui permet de concevoir une application cohérente pour les plateformes mobiles et web. Le système s'appuie sur un backend sécurisé développé avec \textbf{FastAPI}, une couche de persistance basée sur \textbf{SQLAlchemy} et des services dédiés à l'authentification, à la gestion des données, aux notifications et à la modération des contenus. Cette architecture client--serveur garantit la modularité, la maintenabilité et l'évolutivité de la plateforme.

Le présent rapport est structuré en sept chapitres. Le premier chapitre présente le cadre général du projet, l'organisme d'accueil, l'étude de l'existant, la problématique ainsi que la méthodologie de travail adoptée. Le deuxième chapitre est consacré à l'analyse et à la spécification des besoins fonctionnels et non fonctionnels du système. Le troisième chapitre décrit l'environnement technique, les technologies retenues et les choix d'architecture. Les quatrième et cinquième chapitres présentent les réalisations du projet organisées en releases et en sprints. Le sixième chapitre traite l'intégration, les tests et le déploiement de la solution. Enfin, le septième chapitre clôture ce rapport par une conclusion générale et présente les principales perspectives d'évolution de la plateforme \textit{EcoRewind}.
